\newpage
\section{第19章：案例研究——Featherweight Java}

\begin{taplauthorsolutionentrybox}
\phantomsection
\label{sol:author:19.4.1}
\textbf{19.4.1 解答}

每项类声明都必须带有 \texttt{extends} 子句，而这些子句又不得形成环。因此，
从任何类开始沿 \texttt{extends} 链不断向上查找，最终都必定到达
\texttt{Object}。所以无需另设规则，已有规则就能导出每个类都是
\texttt{Object} 的子类型。

\taplsolutionbacklink{ex:19.4.1}{返回练习19.4.1}
\end{taplauthorsolutionentrybox}

\begin{taplauthorsolutionentrybox}
\phantomsection
\label{sol:author:19.4.2}
\textbf{19.4.2 解答}

一个显而易见的改进，是把类型转换的三条类型规则合并成一条：
\[
\inferrule*[right=\textsc{T-Cast}]
 {\Gamma\vdash t_0:D}
 {\Gamma\vdash(C)t_0:C}
\]
并取消“荒谬转换”这一概念。另一项改进是省略构造函数，因为它们本来就不执行
任何实质计算。

\taplsolutionbacklink{ex:19.4.2}{返回练习19.4.2}
\end{taplauthorsolutionentrybox}

\begin{taplauthorsolutionentrybox}
\phantomsection
\label{sol:author:19.4.6}
\textbf{19.4.6 解答}

\begin{enumerate}
\item 在 FJ 中加入接口的形式化定义并不困难。

\item 假设声明以下接口：
\begin{lstlisting}
interface A {}
interface B {}
interface C extends A,B {}
interface D extends A,B {}
\end{lstlisting}
此时 (A) 与 (B) 都是 (C,D) 的共同上界，但二者之间没有最小者，因而
(C,D) 没有最小上界。

\item 不采用条件表达式通常使用的算法式规则
\[
\inferrule
 {\Gamma\vdash t_1:\mathsf{boolean}\\
  \Gamma\vdash t_2:E_2\\
  \Gamma\vdash t_3:E_3}
 {\Gamma\vdash t_1?t_2:t_3:E_2\mathbin{\lor}E_3}
\]
Java 使用以下两条受限规则：
\[
\inferrule
 {\Gamma\vdash t_1:\mathsf{boolean}\\
  \Gamma\vdash t_2:E_2\\
  \Gamma\vdash t_3:E_3\\
  E_2\subtype E_3}
 {\Gamma\vdash t_1?t_2:t_3:E_3}
\]
\[
\inferrule
 {\Gamma\vdash t_1:\mathsf{boolean}\\
  \Gamma\vdash t_2:E_2\\
  \Gamma\vdash t_3:E_3\\
  E_3\subtype E_2}
 {\Gamma\vdash t_1?t_2:t_3:E_2}
\]
这些规则看起来合理，却无法与 FJ 所用的小步操作语义良好配合：类型保持性
实际上不成立，很容易构造反例。
\end{enumerate}

\taplsolutionbacklink{ex:19.4.6}{返回练习19.4.6}
\end{taplauthorsolutionentrybox}

\begin{taplauthorsolutionentrybox}
\phantomsection
\label{sol:author:19.4.7}
\textbf{19.4.7 解答}

有些出人意料的是，处理 \texttt{super} 比处理 \texttt{self} 更困难，因为
解释 \texttt{super} 时，必须记住当前正在执行的方法体最初定义在哪个类中。
至少有两种办法可以保存这项信息：
\begin{enumerate}
\item 给项添加标注，指明应从哪里开始查找其中的 \texttt{super}；
\item 增加一个预处理步骤，重写整个类表：把对 \texttt{super} 的引用改写成
对 \texttt{this} 的引用，同时使用经过改名的方法名来记录它们来自哪个类。
\end{enumerate}

\taplsolutionbacklink{ex:19.4.7}{返回练习19.4.7}
\end{taplauthorsolutionentrybox}

\begin{taplauthorsolutionentrybox}
\phantomsection
\label{sol:author:19.5.1}
\textbf{19.5.1 解答}

证明主定理前，先建立几个辅助引理。与以往一样，关键是把类型判断与替换联系
起来的引理 A.14。

\medskip
\textbf{引理 A.13}
若
\(\mathit{mtype}(m,D)=\overline C\to C_0\)，则对每个
\(C\subtype D\) 都有
\(\mathit{mtype}(m,C)=\overline C\to C_0\)。

\textit{证明。}
对 \(C\subtype D\) 的推导作直接归纳。若
\(CT(C)=\mathsf{class}\ C\ \mathsf{extends}\ E\{\ldots\}\)，无论
\(m\) 是否在 \(CT(C)\) 中定义，
\(\mathit{mtype}(m,C)\) 都应与 \(\mathit{mtype}(m,E)\) 相同。

\medskip
\textbf{引理 A.14（项替换保持类型）}
若
\[
 \Gamma,\overline x:\overline B\vdash t:D,
 \qquad
 \Gamma\vdash\overline s:\overline A,
 \qquad
 \overline A\subtype\overline B,
\]
则存在 \(C\subtype D\)，使
\[
 \Gamma\vdash[\overline x\mapsto\overline s]t:C
\]

\textit{证明。}
对 \(\Gamma,\overline x:\overline B\vdash t:D\) 的推导作归纳。基本思路
与带子类型化 lambda 演算完全相同，细节略有差异；最后两种情形最值得注意。

\textbf{情形 T-Var：}
\(t=x\)，且 \(x:D\in\Gamma,\overline x:\overline B\)。若
\(x\notin\{\overline x\}\)，替换不改变 \(x\)，结论立即成立。若
\(x=x_i\) 且 \(D=B_i\)，则替换结果为 \(s_i\)；由
\(\Gamma\vdash s_i:A_i\) 及 \(A_i\subtype B_i\)，取 \(C=A_i\) 即可。

\textbf{情形 T-Field：}
\(t=t_0.f_i\)，且
\[
 \Gamma,\overline x:\overline B\vdash t_0:D_0,
 \qquad
 \mathit{fields}(D_0)=\overline C\ \overline f,
 \qquad D=C_i
\]
由归纳假设，存在 \(C_0\subtype D_0\)，使
\(\Gamma\vdash[\overline x\mapsto\overline s]t_0:C_0\)。容易验证，对某个
字段类型序列 \(\overline D\)，有
\[
 \mathit{fields}(C_0)=
 \mathit{fields}(D_0),\overline D\ \overline g
\]
所以 \(f_i\) 在 \(C_0\) 中仍有类型 \(C_i\)。由 T-Field 得
\[
 \Gamma\vdash
 ([\overline x\mapsto\overline s]t_0).f_i:C_i
\]

\textbf{情形 T-Invk：}
\(t=t_0.m(\overline t)\)，且
\[
\begin{gathered}
 \Gamma,\overline x:\overline B\vdash t_0:D_0,
 \qquad \mathit{mtype}(m,D_0)=\overline E\to D,\\
 \Gamma,\overline x:\overline B\vdash\overline t:\overline D,
 \qquad \overline D\subtype\overline E
\end{gathered}
\]
由归纳假设，存在 \(C_0,\overline C\)，使
\[
\begin{gathered}
 \Gamma\vdash[\overline x\mapsto\overline s]t_0:C_0,
 \quad C_0\subtype D_0,\\
 \Gamma\vdash[\overline x\mapsto\overline s]\overline t:\overline C,
 \quad \overline C\subtype\overline D
\end{gathered}
\]
由引理 A.13，
\(\mathit{mtype}(m,C_0)=\overline E\to D\)；再由子类型关系的传递性，
\(\overline C\subtype\overline E\)。于是 T-Invk 给出
\[
 \Gamma\vdash
 ([\overline x\mapsto\overline s]t_0).m
 ([\overline x\mapsto\overline s]\overline t):D
\]

\textbf{情形 T-New：}
\(t=\mathsf{new}\ D(\overline t)\)，且
\[
 \mathit{fields}(D)=\overline D\ \overline f,
 \quad
 \Gamma,\overline x:\overline B\vdash\overline t:\overline C,
 \quad
 \overline C\subtype\overline D
\]
归纳假设给出某个 \(\overline E\subtype\overline C\)，满足
\(\Gamma\vdash[\overline x\mapsto\overline s]\overline t:\overline E\)。
由传递性，\(\overline E\subtype\overline D\)，再用 T-New 即得结论。

\textbf{情形 T-UCast：}
\(t=(D)t_0\)，且
\(\Gamma,\overline x:\overline B\vdash t_0:C\)、\(C\subtype D\)。归纳
假设给出 \(E\subtype C\)，满足
\(\Gamma\vdash[\overline x\mapsto\overline s]t_0:E\)。由传递性
\(E\subtype D\)，再用 T-UCast 得
\(\Gamma\vdash(D)([\overline x\mapsto\overline s]t_0):D\)。

\textbf{情形 T-DCast：}
\(t=(D)t_0\)，且
\(\Gamma,\overline x:\overline B\vdash t_0:C\)、\(D\subtype C\)、
\(D\ne C\)。归纳假设给出 \(E\subtype C\)，满足
\(\Gamma\vdash[\overline x\mapsto\overline s]t_0:E\)。若
\(E\subtype D\) 或 \(D\subtype E\)，分别用 T-UCast 或 T-DCast 得到结果；
若 \(D\nsubtype E\) 且 \(E\nsubtype D\)，则由 T-SCast 得到结果，并产生一条
荒谬转换警告。

\textbf{情形 T-SCast：}
\(t=(D)t_0\)，且
\(\Gamma,\overline x:\overline B\vdash t_0:C\)、\(D\nsubtype C\)、
\(C\nsubtype D\)。归纳假设给出 \(E\subtype C\)，满足
\(\Gamma\vdash[\overline x\mapsto\overline s]t_0:E\)。由 FJ 中每个类只有
一个直接超类可知 \(E\nsubtype D\)：否则若同时有 \(E\subtype C\) 和
\(E\subtype D\)，则 \(C\subtype D\) 与 \(D\subtype C\) 中必有一个成立，
与前提矛盾。因此可用 T-SCast 得到结果，并产生荒谬转换警告。

\medskip
\textbf{引理 A.15（弱化）}
若 \(\Gamma\vdash t:C\)，则 \(\Gamma,x:D\vdash t:C\)。

\textit{证明。} 对类型推导作直接归纳。

\medskip
\textbf{引理 A.16}
若
\[
 \mathit{mtype}(m,C_0)=\overline D\to D,
 \qquad
 \mathit{mbody}(m,C_0)=(\overline x,t),
\]
则存在 \(D_0\) 及 \(C\subtype D\)，使
\[
 C_0\subtype D_0,
 \qquad
 \overline x:\overline D,\mathsf{this}:D_0\vdash t:C
\]

\textit{证明。}
对 \(\mathit{mbody}(m,C_0)\) 的推导作归纳。基例中 \(m\) 就定义在
\(C_0\) 内；类表良构保证 T-Method 已经导出
\(\overline x:\overline D,\mathsf{this}:C_0\vdash t:C\)。归纳步骤同样
直接。

现在证明类型安全性定理。对 \(t\trans t'\) 的推导作归纳，并按最后使用的规则
分类。倒数第二个 T-DCast 子情形会产生荒谬转换警告。

\textbf{情形 E-ProjNew：}
设
\[
 t=(\mathsf{new}\ C_0(\overline v)).f_i,
 \qquad t'=v_i,
 \qquad \mathit{fields}(C_0)=\overline D\ \overline f
\]
从 \(t\) 的形式可知，\(\Gamma\vdash t:C\) 的最后一条规则必为 T-Field，
其前提是
\(\Gamma\vdash\mathsf{new}\ C_0(\overline v):D_0\)，且
\(C=D_i\)。后一判断的最后一条规则又必为 T-New，其前提包括
\(\Gamma\vdash\overline v:\overline C\) 与
\(\overline C\subtype\overline D\)，并且 \(D_0=C_0\)。特别地，
\(\Gamma\vdash v_i:C_i\)，而 \(C_i\subtype D_i=C\)，结论成立。

\textbf{情形 E-InvkNew：}
设
\[
\begin{gathered}
 t=(\mathsf{new}\ C_0(\overline v)).m(\overline u),\\
 t'=[\overline x\mapsto\overline u,
     \mathsf{this}\mapsto\mathsf{new}\ C_0(\overline v)]t_0,\\
 \mathit{mbody}(m,C_0)=(\overline x,t_0)
\end{gathered}
\]
\(\Gamma\vdash t:C\) 的最后两条规则必为 T-Invk 与 T-New，其前提包括
\[
 \Gamma\vdash\mathsf{new}\ C_0(\overline v):C_0,
 \quad
 \Gamma\vdash\overline u:\overline C,
 \quad
 \overline C\subtype\overline D,
 \quad
 \mathit{mtype}(m,C_0)=\overline D\to C
\]
由引理 A.16，存在 \(D_0\) 与 \(B\subtype C\)，满足
\(C_0\subtype D_0\) 且
\(\overline x:\overline D,\mathsf{this}:D_0\vdash t_0:B\)。由引理 A.15
可把 \(\Gamma\) 加入环境；再用引理 A.14 同时替换形参与
\texttt{this}，得到某个 \(E\subtype B\)，使
\(\Gamma\vdash t':E\)。由传递性 \(E\subtype C\)，完成此情形。

\textbf{情形 E-CastNew：}
设
\[
 t=(D)(\mathsf{new}\ C_0(\overline v)),
 \qquad C_0\subtype D,
 \qquad t'=\mathsf{new}\ C_0(\overline v)
\]
\(\Gamma\vdash t:C\) 的推导只能以 T-UCast 结束；若以 T-SCast 或
T-DCast 结束，就会与 \(C_0\subtype D\) 矛盾。T-UCast 的前提给出
\(\Gamma\vdash\mathsf{new}\ C_0(\overline v):C_0\) 与 \(D=C\)，所以取
\(C'=C_0\) 即可。

同余规则的情形都很直接，只展示 E-Cast。设
\(t=(D)t_0\)、\(t'=(D)t_0'\)、\(t_0\trans t_0'\)。根据原推导最后使用
的类型规则，分三种子情形。

\textbf{子情形 T-UCast：}
已知 \(\Gamma\vdash t_0:C_0\)、\(C_0\subtype D\)、\(D=C\)。归纳假设
给出某个 \(C_0'\subtype C_0\)，使 \(\Gamma\vdash t_0':C_0'\)。由传递性
\(C_0'\subtype C\)，所以 T-UCast 导出
\(\Gamma\vdash(C)t_0':C\)，且不增加荒谬转换警告。

\textbf{子情形 T-DCast：}
已知 \(\Gamma\vdash t_0:C_0\)、\(D\subtype C_0\)、\(D=C\)。归纳假设
给出 \(C_0'\subtype C_0\)，使 \(\Gamma\vdash t_0':C_0'\)。若
\(C_0'\subtype C\) 或 \(C\subtype C_0'\)，分别用 T-UCast 或 T-DCast，
不增加荒谬转换警告；若两边都不是子类型，则用 T-SCast，并产生一条荒谬转换
警告。

\textbf{子情形 T-SCast：}
已知 \(\Gamma\vdash t_0:C_0\)、\(D\nsubtype C_0\)、\(C_0\nsubtype D\)、
\(D=C\)。归纳假设给出 \(C_0'\subtype C_0\)，使
\(\Gamma\vdash t_0':C_0'\)。于是仍有
\(C_0'\nsubtype C\) 与 \(C\nsubtype C_0'\)，所以 T-SCast 导出
\(\Gamma\vdash(C)t_0':C\)，并保留荒谬转换警告。

\taplsolutionbacklink{thm:19.5.1}{返回定理19.5.1}
\end{taplauthorsolutionentrybox}
